% Einheit 4 von 4 – Periodische Vorgänge modellieren (bank/trigonometrische-funktionen/e4.jsonl, zusammenbau v0.1)
%% TODO Zweigzeile Teil 1: was hier gelernt wird (Satz in Schülersprache aus dem Einheitstitel) – Einheit 4
\einheitenkopf[e4]{Einheit 4 von 4 · Periodische Vorgänge modellieren}
\zweigzeile{ab Kl. 10 · keine P10-Aufgabe}
\setcounter{aufgabe}{20}

%% TODO Titel: Ich-kann-Satz fehlt in der Bank (Platzhalter: Kettenname) – Nr. 21
%% TODO Anweisung: ein Satz über der ersten Teilaufgabe fehlt in der Bank – Nr. 21
\begin{aufgabe}{Periodische Vorgänge modellieren}
\begin{teile}
\teil Wiederholt sich das? Die Höhe des Ventils am Fahrradreifen bei gleichmäßiger Fahrt.\\ \kreuz{wiederholt sich}\kreuz{wiederholt sich nicht}
\end{teile}
%% TODO Darstellung für schwach fehlt in der Bank (trigonometrische-funktionen-e4-k1-s1-v1; Abschnitt „Für schwache Schüler“)
\swz[3]{Der Ausschlag eines Uhrpendels nach links und rechts – periodisch oder nicht? Begründe in einem Satz.}{}
%% TODO Darstellung für schwach fehlt in der Bank (trigonometrische-funktionen-e4-k1-s1-v2; Abschnitt „Für schwache Schüler“)
\swz[3]{Ein Kontostand, auf den jeden Monat derselbe Betrag eingezahlt wird – periodisch oder nicht? Begründe in einem Satz.}{}
%% TODO Darstellung für schwach fehlt in der Bank (trigonometrische-funktionen-e4-k1-s1-v3; Abschnitt „Für schwache Schüler“)
\swz[3]{Die Höhe der Spitze des Minutenzeigers einer Uhr – periodisch oder nicht? Begründe in einem Satz.}{}
%% TODO Darstellung für schwach fehlt in der Bank (trigonometrische-funktionen-e4-k1-s1-v4; Abschnitt „Für schwache Schüler“)
\swz[3]{Der Anhalteweg eines Autos bei immer höherer Geschwindigkeit – periodisch oder nicht? Begründe in einem Satz.}{}
%% TODO Darstellung für schwach fehlt in der Bank (trigonometrische-funktionen-e4-k1-s1-v5; Abschnitt „Für schwache Schüler“)
\swz[3]{Die Höhe einer Gondel im Riesenrad bei gleichmäßiger Fahrt – periodisch oder nicht? Begründe in einem Satz.}{}
\swfrage{Was bleibt gleich, was ändert sich?}
%% TODO Darstellung für schwach fehlt in der Bank (trigonometrische-funktionen-e4-k1-s2-v1; Abschnitt „Für schwache Schüler“)
\swz{Ein Leuchtfeuer ist 2 s hell, dann 4 s dunkel, dann wieder hell. Wie lang ist die Periode?}{}
\swz{Temperatur in einem Gewächshaus – Größtwert und Kleinstwert?}{\sachtabelle{lcccccc}{Uhrzeit & 2 & 6 & 10 & 14 & 18 & 22}{Temperatur in °C & 15 & 13 & 17 & 25 & 27 & 21}}
%% TODO Darstellung für schwach fehlt in der Bank (trigonometrische-funktionen-e4-k1-s4-v1; Abschnitt „Für schwache Schüler“)
\swz{Größtwert 27 °C, Kleinstwert 13 °C – Mittellinie und Amplitude?}{}
%% TODO Darstellung für schwach fehlt in der Bank (trigonometrische-funktionen-e4-k1-s5-v1; Abschnitt „Für schwache Schüler“)
\swz{Mittellinie 20 °C, Amplitude 7 °C, Periode 24 h – wie lautet die Gleichung y = a · sin(b · t) + d?}{}
\begin{teile}
\teil Eine Gondel fährt zwischen 2 m und 34 m Höhe, eine Runde dauert 8 min. Welche Gleichung passt?\\ \kreuz{y = 16 · sin(45 · t) + 18}\\ \kreuz{y = 18 · sin(45 · t) + 16}\\ \kreuz{y = 16 · sin(8 · t) + 18}
\end{teile}
%% TODO Darstellung für schwach fehlt in der Bank (trigonometrische-funktionen-e4-k1-s7-v1; Abschnitt „Für schwache Schüler“)
\swz[3]{Für die Temperatur gilt y = 7 · sin(15 · t) + 20 (t in Stunden nach 8 Uhr, y in °C). Berechne y für t = 4 und deute den Wert.}{}
%% TODO Darstellung für schwach fehlt in der Bank (trigonometrische-funktionen-e4-k1-s8-v1; Abschnitt „Für schwache Schüler“)
\swz{y = 7 · sin(15 · t) + 20 (t in Stunden nach 8 Uhr) – wann ist es am wärmsten?}{}
\swz{Lies mit Hilfslinien ab: Wie hoch ist die Gondel nach 1 Minute?}{\begin{ksys}[xmin=0,xmax=16,ymin=0,ymax=36,xstep=1,ystep=2,xlabel=t in min,ylabel=Höhe in m,ablesen]\funktion{16*sin(45*\x)+18}{f}\end{ksys}}
%% TODO Darstellung für schwach fehlt in der Bank (trigonometrische-funktionen-e4-k1-s10-v1; Abschnitt „Für schwache Schüler“)
\swz[3]{Das Modell y = 7 · sin(15 · t) + 20 beschreibt einen Sommertag im Gewächshaus. Beurteile, ob es für eine ganze Woche passt, und nenne eine Grenze.}{}
\swa{Fülle die Tabelle aus: Welcher Funktionstyp wiederholt sich, welcher nimmt je Schritt immer gleich viel zu?}{\sachtabelle{lcccc}{ & Gerade & Parabel & Exponentialkurve & Sinuskurve}{wiederholt sich & \leerzelle & \leerzelle & \leerzelle & \leerzelle \\ gleiche Zunahme je Schritt & \leerzelle & \leerzelle & \leerzelle & \leerzelle}}
\begin{teile}
\teil Ein Guthaben von 500 € wächst jedes Jahr um 3 \%. Welcher Graph passt? Kreuze an und begründe für die beiden anderen, warum sie nicht passen. \hfill (P10 2019 OS)\\ \kreuz{A}\kreuz{B}\kreuz{C}
\end{teile}
\swz[3]{Die Tageslänge in einer Stadt schwankt zwischen 8 h und 16,4 h, die Periode ist 365 Tage; t zählt die Tage ab dem 21. März. a) Welche Gleichung passt: (1) y = 4,2 · sin(0,99 · t) + 12,2, (2) y = 12,2 · sin(0,99 · t) + 4,2 oder (3) y = 4,2 · sin(365 · t) + 12,2? b) Welcher Graph, A oder B, passt? c) Berechne y für t = 100 und deute den Wert im Sachzusammenhang.}{\begin{ksys}[xmin=0,xmax=360,ymin=0,ymax=22,xstep=60,ystep=2,xlabel=t in Tagen,ylabel=Stunden,ablesen]\funktion{4.2*sin(0.99*\x)+12.2}{A}\funktion{8.4*sin(0.99*\x)+12.2}{B}\end{ksys}}
\end{aufgabe}

%% TODO Titel: Ich-kann-Satz fehlt in der Bank (Platzhalter: Kettenname) – Nr. 22
%% TODO Anweisung: ein Satz über der ersten Teilaufgabe fehlt in der Bank – Nr. 22
\begin{aufgabe}{Periodische Vorgänge modellieren – Fehler finden, Begründen}
\swz[3]{Zum Graphen der Höhe einer Riesenradgondel über der Zeit sagt Emma: Der Graph zeigt den Weg der Gondel – sie fährt eine Welle. Finde den Fehler.}{}
\swfrage{Begründe, warum ein Guthaben mit festem Zinssatz keine Sinuskurve ergibt.}
\end{aufgabe}

\uebersichtskasten{Periodisch heißt: Der Vorgang wiederholt sich immer nach derselben Zeitspanne – Tageslänge im Jahreslauf, Ebbe und Flut, Herzschlag, Ton und Welle. Diese Zeitspanne ist die Periode. \\ Vom Vorgang zur Gleichung: Mittellinie d ist der Wert genau zwischen Größt- und Kleinstwert, Amplitude a ist ihr halber Abstand, und aus der Periode ergibt sich b als Vollkreis geteilt durch die Periode. \\ Wassertiefe zwischen 2 m und 6 m, Wiederholung alle 12 Stunden: Mittellinie 4, Amplitude 2, Periode 12. \\ Modell und Wirklichkeit: Ein echter Vorgang ist meist nur annähernd periodisch – die Rechnung liefert Näherungswerte, keine genauen Messwerte. \\ Welcher Funktionstyp passt? Gerade: immer gleiche Zunahme. Parabel: eine Spitze, dann zurück. Exponentialkurve: gleicher Faktor je Schritt, wird immer steiler oder immer flacher. Sinuskurve: wiederholt sich.}
